\documentclass[11pt,a4paper]{article}
\usepackage[margin=24mm]{geometry}
\usepackage{amsmath,amssymb,bm,booktabs,array,microtype}
\usepackage[T1]{fontenc}
\usepackage{lmodern}
\usepackage[hidelinks]{hyperref}
\usepackage{xcolor}

\definecolor{rileyblue}{RGB}{35,83,120}
\hypersetup{pdftitle={Single-Element Quadratic Saddle Deformation for a Quad9 Element}}

\newcommand{\vect}[1]{\bm{#1}}
\newcommand{\dd}{\mathrm{d}}

\title{\textbf{Single-Element Quadratic Saddle Deformation}\\
\large Manufactured displacement field for a nine-node quadrilateral (Quad9) element}
\author{}
\date{}

\begin{document}
\maketitle

\section{Purpose}

This note defines a compact, purely quadratic displacement field for a single Quad9 finite element. The field is intended for controlled image-rendering and digital image correlation (DIC) studies in which the deformation amplitude is ramped from zero to one pixel while the geometry, texture, camera and image-processing settings remain fixed.

The field provides a useful progression beyond rigid and affine motion because it contains no constant or linear displacement terms, exercises both image-coordinate displacement components, and is represented exactly by the biquadratic interpolation of a Quad9 element.

\section{Element coordinates}

Let the element occupy a square plate-centred domain
\begin{equation}
  x\in[x_c-L_x,\,x_c+L_x],
  \qquad
  y\in[y_c-L_y,\,y_c+L_y].
\end{equation}
Define the normalised element coordinates
\begin{equation}
  \xi=\frac{x-x_c}{L_x},
  \qquad
  \eta=\frac{y-y_c}{L_y},
  \qquad
  (\xi,\eta)\in[-1,1]^2.
\end{equation}
The shared single-element generator uses a plate side length $P=260$ and an
ROI side length $R=256$.  Therefore the saddle element uses
\begin{equation}
  x_c=y_c=0,
  \qquad
  L_x=L_y=P/2=130.
\end{equation}
The generator's $256\times256$ reference camera gives a physical pixel width
$\Delta=R/256=1$.  The ROI is used for UV mapping and padded texture coverage;
it does not change the plate-coordinate normalisation above.

\section{Quadratic saddle displacement field}

Let $d$ denote the prescribed maximum displacement amplitude in physical units
(one unit is one reference-camera pixel). Define
\begin{equation}
  \boxed{
  \vect{u}(\xi,\eta;d)
  =
  \begin{bmatrix}
    u_x\\[2pt]
    u_y
  \end{bmatrix}
  =
  \frac{d}{2}
  \begin{bmatrix}
    \xi^2-\eta^2\\[2pt]
    2\xi\eta
  \end{bmatrix}.}
  \label{eq:saddle}
\end{equation}
Equivalently,
\begin{equation}
  u_x(\xi,\eta;d)=\frac{d}{2}(\xi^2-\eta^2),
  \qquad
  u_y(\xi,\eta;d)=d\xi\eta.
\end{equation}
The amplitude can be swept as
\begin{equation}
  d\in\{0,0.1,0.2,\ldots,1.0\}\ \text{pixels}.
\end{equation}
This creates a one-parameter family of geometrically similar deformation fields, with the displacement and displacement gradient scaling linearly with $d$.

\section{Amplitude normalisation}

The squared displacement magnitude is
\begin{align}
  \|\vect{u}\|^2
  &=\frac{d^2}{4}(\xi^2-\eta^2)^2+d^2\xi^2\eta^2\\
  &=\frac{d^2}{4}(\xi^2+\eta^2)^2.
\end{align}
Hence
\begin{equation}
  \boxed{\|\vect{u}(\xi,\eta;d)\|
  =\frac{d}{2}(\xi^2+\eta^2).}
\end{equation}
Since $\xi^2+\eta^2\le2$ on the square element,
\begin{equation}
  \max_{[-1,1]^2}\|\vect{u}\|=d,
\end{equation}
with the maximum attained at all four corner nodes. Therefore, setting $d=1$ gives an exact maximum nodal displacement of one pixel.

The displacement vanishes at the element centre and changes direction around it. Along the horizontal and vertical centre lines, the motion lies entirely in the $x$ direction; along the diagonals, the motion lies entirely in the $y$ direction.

\section{Quad9 interpolation}

The one-dimensional quadratic Lagrange basis functions on $\{-1,0,1\}$ are
\begin{equation}
  \ell_{-1}(s)=\frac{1}{2}s(s-1),
  \qquad
  \ell_0(s)=1-s^2,
  \qquad
  \ell_{+1}(s)=\frac{1}{2}s(s+1).
\end{equation}
The nine tensor-product Quad9 shape functions are
\begin{equation}
  N_{ab}(\xi,\eta)=\ell_a(\xi)\ell_b(\eta),
  \qquad a,b\in\{-1,0,1\}.
\end{equation}
The finite-element displacement is
\begin{equation}
  \vect{u}_h(\xi,\eta)
  =\sum_{a,b\in\{-1,0,1\}}
  N_{ab}(\xi,\eta)\vect{u}_{ab},
\end{equation}
where $\vect{u}_{ab}$ is obtained by evaluating Eq.~\eqref{eq:saddle} at the corresponding node.

Because $u_x$ and $u_y$ are polynomials contained in the biquadratic space
\begin{equation}
  \mathcal{Q}_2
  =\operatorname{span}\{\xi^i\eta^j:0\le i,j\le2\},
\end{equation}
the Quad9 interpolation reproduces the manufactured field exactly:
\begin{equation}
  \boxed{\vect{u}_h(\xi,\eta)=\vect{u}(\xi,\eta)
  \quad\text{for all }(\xi,\eta)\in[-1,1]^2.}
\end{equation}
There is therefore no finite-element interpolation error in the imposed displacement field. Any observed image error arises from the rendering, texture representation, pixel integration, image quantisation or subsequent measurement process.

\section{Quad9 nodal displacement values}

Using the conventional node layout
\[
\begin{array}{ccc}
4&7&3\\
8&9&6\\
1&5&2
\end{array}
\]
with the natural coordinates shown in Table~\ref{tab:nodes}, the exact nodal displacements are as follows.

\begin{table}[ht]
\centering
\caption{Quad9 nodal values for the quadratic saddle field.}
\label{tab:nodes}
\renewcommand{\arraystretch}{1.2}
\begin{tabular}{ccccc}
\toprule
Node & $\xi$ & $\eta$ & $u_x$ & $u_y$\\
\midrule
1 & $-1$ & $-1$ & $0$ & $+d$\\
2 & $+1$ & $-1$ & $0$ & $-d$\\
3 & $+1$ & $+1$ & $0$ & $+d$\\
4 & $-1$ & $+1$ & $0$ & $-d$\\
5 & $0$  & $-1$ & $-d/2$ & $0$\\
6 & $+1$ & $0$  & $+d/2$ & $0$\\
7 & $0$  & $+1$ & $-d/2$ & $0$\\
8 & $-1$ & $0$  & $+d/2$ & $0$\\
9 & $0$  & $0$  & $0$ & $0$\\
\bottomrule
\end{tabular}
\end{table}

\section{Displacement gradient and small-strain field}

The natural-coordinate derivatives are
\begin{equation}
  \frac{\partial u_x}{\partial\xi}=d\xi,
  \qquad
  \frac{\partial u_x}{\partial\eta}=-d\eta,
  \qquad
  \frac{\partial u_y}{\partial\xi}=d\eta,
  \qquad
  \frac{\partial u_y}{\partial\eta}=d\xi.
\end{equation}
Using
\begin{equation}
  \frac{\partial}{\partial x}=\frac{1}{L_x}\frac{\partial}{\partial\xi},
  \qquad
  \frac{\partial}{\partial y}=\frac{1}{L_y}\frac{\partial}{\partial\eta},
\end{equation}
the image-space displacement gradient is
\begin{equation}
  \boxed{
  \nabla\vect{u}
  =
  \begin{bmatrix}
    d\xi/L_x & -d\eta/L_y\\[4pt]
    d\eta/L_x & d\xi/L_y
  \end{bmatrix}.}
\end{equation}
For a square element, $L_x=L_y=L$, this becomes
\begin{equation}
  \nabla\vect{u}
  =\frac{d}{L}
  \begin{bmatrix}
    \xi & -\eta\\
    \eta & \xi
  \end{bmatrix}.
\end{equation}
The infinitesimal strain tensor is
\begin{equation}
  \vect{\varepsilon}
  =\frac{1}{2}\left(\nabla\vect{u}+\nabla\vect{u}^{\mathsf T}\right).
\end{equation}
For the square case,
\begin{equation}
  \boxed{
  \vect{\varepsilon}
  =\frac{d\xi}{L}
  \begin{bmatrix}
    1&0\\
    0&1
  \end{bmatrix}.}
\end{equation}
Thus the symmetric part is a linearly varying equibiaxial strain field, while the antisymmetric part is a linearly varying local rotation proportional to $\eta$. Although the displacement field has a saddle-like directional pattern, its first-order strain and rotation components remain simple and interpretable.

For $L=130$ and $d=1$, the characteristic gradient scale is
\begin{equation}
  \frac{d}{L}=\frac{1}{130}\approx7.6923\times10^{-3},
\end{equation}
which exercises quadratic mapping while remaining consistent with the small
one-pixel displacement scale of the rigid and affine cases.

\section{Recommended parameter sweep}

A minimal study can use
\begin{equation}
  d_k=0.1k\ \text{pixels},
  \qquad k=0,1,\ldots,10.
\end{equation}
At each amplitude, keep all other quantities fixed and compare:
\begin{itemize}
  \item the analytic or independently generated image;
  \item Riley renders across texture oversampling levels;
  \item Riley renders across pixel-box supersampling levels;
  \item floating-point image error;
  \item digitised image error at the selected camera bit depths;
  \item optionally, a fixed-setting DIC measurement against the prescribed field.
\end{itemize}
The field complements rigid and affine tests by providing the sequence
\begin{equation}
  \text{constant displacement}
  \;\longrightarrow\;
  \text{linear displacement}
  \;\longrightarrow\;
  \text{quadratic displacement}.
\end{equation}

\section{Implementation summary}

For each amplitude $d$:
\begin{enumerate}
  \item create a single planar Quad9 element over the centred $260\times260$ plate;
  \item assign the nodal displacements from Table~\ref{tab:nodes};
  \item use the $256\times256$ ROI and any texture pad only for rendering and texture support;
  \item render the undeformed and deformed states using identical camera and texture settings;
  \item crop or analyse only the central active region when evaluating convergence and DIC error.
\end{enumerate}

The manufactured field is exactly representable by the Quad9 element, has a maximum displacement equal to the scalar ramp parameter $d$, and provides a compact non-affine deformation case for the initial convergence experiments.

\end{document}
